\documentclass[../../main.tex]{subfiles}
\begin{document}

\paragraph{【乙】}
$n$枚のカードを積んだ山があり，
各カードには上から順番に1から$n$まで番号がつけられている．
ただし$n\geqq2$とする．
このカードの山に対して次の試行を繰り返す．
1回の試行では，
一番上のカードを取り，
山の一番上にもどすか，
あるいはいずれかのカードの下に入れるという操作を行う．
これら$n$通りの操作はすべて同じ確率であるとする．
$n$回の試行を終えたとき，
最初一番下にあったカード(番号$n$)が山の一番上にきている確率を求めよ．

\paragraph{【甲】}
$x$，$y$は$x \neq 1$，$y \neq 1$をみたす正の数で，不等式
\begin{align}
  \log_x y + \log_y x > 2 + (\log_x 2)(\log_y 2)
\end{align}
をみたすとする．このとき$x$，$y$の組$(x,y)$の範囲を座標平面上に図示せよ．

\end{document}
